\section{Toy Generation and Validation Studies}
\label{sec:toys}

\begin{figure}[H]
\centering
\graphicorplaceholder{0.98\linewidth}{guard_band_figs/hps_gpr_bump_hunt_flowchart.png}
\caption{HPS GPR bump-hunt analysis and closure-validation flow. The shared
per-mass setup defines the spectrum, signal template, blind/extraction window,
sideband-training region, GP refit, and amplitude extraction. The nominal-data path
converts the profiled yield upper limit to $\eps^2_{\mathrm{UL}}$, while the
functional-form path injects known signals into smooth-background toys and tests signal
recovery, pull calibration, $Z$ calibration, uncertainty consistency, and leakage into
the sideband-training region.}
\label{fig:hps-gpr-flowchart-section5}
\end{figure}

\subsection{Validation of the Statistical Framework}

The closure tests in this section validate the local extraction and are not themselves
confidence-level statements. Some historical reach proxies use the one-sided 95\%
Gaussian coefficient, while the current physics-facing result in
Section~\ref{sec:results} is a 90\% CL upper limit obtained with the \CLs{} construction.
The staged suites cover the fully unblinded 2015 sample, the 2016 10\% development
sample, and the regenerated 2021 1\% validation sample. The headline functional-form
figures use full GP refits and, for 2021, the corrected matched-reference widened-window
convention. Superseded and duplicated closure displays are collected in
Appendix~\ref{app:prior-validation-results}.

The validation material is organized around four checks of analysis robustness:

\begin{itemize}
    \item First, how far from the signal peak should the GP sideband training data begin?

    \item  Second, after that training exclusion is fixed, should the $1.64$--$2.25\,\sigma_m$ signal tails be omitted as a guard band or included in the extraction likelihood through a wider blind window?

    \item Third, do the reference-matched full-refit toys close for the 2015, 2016, and 2021 datasets?

    \item   Finally, do the combined pseudoexperiments accumulate sensitivity as expected when statistically independent datasets share one coupling parameter?

\end{itemize}



\subsubsection*{Terms used in the validation plots}

\begin{description}
\item[train-$k$ or $k_{train}$] The GP sideband fit excludes bins with
$|m-m_0|<k\,\sigma_m$ around the tested mass. The selected production value is
train-$2.25$, matched to the $2.25\,\sigma_m$ blind/extraction window.
\item[refmatched] The injected amplitude is defined with the same
background-only refit reference scale, blind geometry, kernel policy, and refit
prescription used in the extraction study:
$A_{\mathrm{inj}}=Z_{\mathrm{inj}}\sigma_{A,\mathrm{ref}}$.
\item[$Z_{\mathrm{inj}}$] The nominal injected signal strength in units of
$\sigma_{A,\mathrm{ref}}$.
\item[$\hat Z$] The extracted signal significance, usually summarized as
$\hat A/\sigma_A$ for toy-closure plots.
\item[pull mean] The ensemble mean of
$(\hat A-A_{\mathrm{inj}})/\sigma_A$. A value near zero indicates little bias.
\item[pull width] The ensemble standard deviation of the pull. A value near one
indicates that the fitted $\sigma_A$ is correctly calibrated.
\item[$\Delta Z$] The significance residual $\hat Z-Z_{\mathrm{inj}}$.
\item[Gaussian-equivalent expected reach proxy] The background-only proxy
$1.64\,\sigma_{A,\mathrm{ref}}/A_{\eps^2=1}$ used to compare window choices. It is a
compact expected-95\% scale, not the final \CLs{} limit.
\end{description}

Throughout this section, the zero-injection row has a special interpretation. At
$Z_{\mathrm{inj}}=0$ no signal is added, so the pull mean, pull width, and
$\Delta Z=\hat Z-Z_{\mathrm{inj}}$ form a spurious-signal test. Centered pulls and
near-zero $\Delta Z$ mean that the GP-plus-template extraction is not manufacturing a
signal in a smooth background-only spectrum. Nonzero injected rows then test whether a
real signal is recovered after the same sideband-refit prescription has been applied.

\subsection{Training exclusion and blind-window width}
\label{subsec:guard-band-construction}

The first decision is the training exclusion: how far from the test mass the sideband bins used by the GP should begin. Figure~\ref{fig:blind225-cartoon} illustrates the adopted wide-blind geometry. For a large injected signal, Gaussian tails remain visible outside the central peak. The training-exclusion studies below show why training directly next to a narrow $1.64\,\sigma_m$ extraction window is unattractive: those tails can be
offered to the GP as ordinary smooth-background information and then be partly absorbed
by the refitted background.

\begin{figure}[H]
\centering
\graphicorplaceholder{0.90\linewidth}{guard_band_figs/guard_band_5sigma_fake_data_gp_band_no_guard_2p25_extended2sigma.png}
\caption{Illustration of the selected $2.25\,\sigma_m$ blind-window construction for a
large injected signal. The GP is trained only outside the widened blinded region, so the
near tails of a Gaussian resonance are not included in the sideband-training data. The
same blinded region is then used by the local signal-extraction likelihood.}
\label{fig:blind225-cartoon}
\end{figure}

The training-exclusion scan starts from the older $1.64\,\sigma_m$ blind window and
varies the outer boundary of the bins excluded from GP training. In this notation,
train-$k$ means that bins with $|m-m_0|<k\,\sigma_m$ are withheld from the GP sideband
fit. The train-$2.25$ choice is not the point with the strongest raw background-only
reach proxy; it is selected from the combined leakage, recovery, and calibration
evidence below.


In the comparison of different blind windows shown in Figure~\ref{fig:guard-band-reach-uncertainty}, train-$1.96$ gives the lowest median Gaussian-equivalent $\eps^2_{95}$ proxy, but it also leaves a much larger fraction of the injected signal template in the GP training region, leading to higher leakage rate and lower signal sensitivity. The train-$2.25$ exclusion is the balanced candidate: it sharply reduces train-region leakage, is more sensitive to injected signal than $1.96$, and retains better expected reach than the wider $2.58$ and $3.0$ exclusions. The compact evidence is summarized in Table~\ref{tab:guard-band-choice-summary}; for the train-$2.25$ choice, the mean signal recovery is about $0.961$ at $Z_{\mathrm{inj}}=3$ and $0.943$ at $Z_{\mathrm{inj}}=5$, with a maximum injected training fraction of about $0.0258$.

\begin{table}[H]
\centering
\begin{tabular}{ccccc}
\toprule
$k_{\mathrm{train}}$ & median expected $\eps^2_{95}$ &
$\langle\hat A/A_{\mathrm{inj}}\rangle_{Z=3}$ &
$\langle\hat A/A_{\mathrm{inj}}\rangle_{Z=5}$ &
max $f_{\mathrm{train}}$ \\
\midrule
1.96 & \num{7.92e-5} & 0.730 & 0.715 & 0.0530 \\
2.25 & \num{9.53e-5} & 0.961 & 0.943 & 0.0258 \\
2.58 & \num{1.08e-4} & 1.052 & 1.036 & 0.0106 \\
3.00 & \num{1.19e-4} & 1.049 & 1.037 & 0.00288 \\
\bottomrule
\end{tabular}
\caption{Training-exclusion comparison summary from the refmatched 2015 closure toys using the older $1.64\,\sigma_m$ blind window. The expected reach column is the background-only Gaussian-equivalent proxy $1.64\,\sigma_{A,\mathrm{ref}}/A_{\eps^2=1}$ defined in the terms list, not the final $CL_s$ limit. The train-$2.25$ point trades a modest reach cost relative to $1.96$ for much smaller training leakage and substantially improved injected-signal recovery.} \label{tab:guard-band-choice-summary}
\end{table}

The figures in this subsection are meant to be read sequentially because each one removes a different ambiguity in the guard-band choice. Figure~\ref{fig:blind225-cartoon} first defines the selected widened blind region and shows the practical reason for the choice: a large Gaussian resonance still has non-negligible tails outside a narrow $1.64\,\sigma_m$ extraction window, so training directly beside that window can feed signal-like counts to the GP sideband fit. Table~\ref{tab:guard-band-choice-summary} then gives the compact numerical trade-off between expected reach, recovered injected signal, and the maximum signal-template fraction left in the training region.

The remaining guard-band figures separate the ingredients behind that table. Figure~\ref{fig:guard-band-reach-uncertainty} shows the cost side of widening the training exclusion: as nearby sideband bins are removed, $\sigma_{A,\mathrm{ref}}$ grows and the Gaussian-equivalent expected reach softens. Figure~\ref{fig:guard-band-leakage-reference} shows the protection side: train-$2.25$ greatly reduces the signal fraction available to the GP training set while preserving the matched reference scale used to define the injected amplitudes. Figure~\ref{fig:guard-band-closure-diagnostics} checks that this protection does not create a new calibration problem by comparing recovery, $\Delta Z$, pull means, and pull widths over the mass and injection grid. Figure~\ref{fig:guard-band-recovery-boxes} then shows the same recovery behavior at the toy-distribution level, making the narrow-guard under-recovery visible beyond the ensemble mean curves. Finally, Figures~\ref{fig:refmatched-blind-geometry-reach}--\ref{fig:refmatched-blind-geometry-deltaz} compare the two blind-window geometries after train-$2.25$ is fixed: expected reach favors including the $1.64$--$2.25\,\sigma_m$ tails in the likelihood, while the background-only, pull-summary, amplitude-recovery, and $\Delta Z$ panels show that the widened blind window keeps the same closure behavior as the explicit guard-band construction.

\begin{figure}[H]
\centering
\begin{subfigure}[t]{0.82\linewidth}
  \centering
  \graphicorplaceholder{\linewidth}{guard_band_figs/expected95_reach_z0.png}
  \caption{Background-only expected $\eps^2_{95}$ proxy.}
\end{subfigure}

\vspace{0.8em}

\begin{subfigure}[t]{0.82\linewidth}
  \centering
  \graphicorplaceholder{\linewidth}{guard_band_figs/uncertainty_vs_reach_z0.png}
  \caption{Reference uncertainty and reach relative to train-$1.96$.}
\end{subfigure}
\caption{Expected-reach cost of widening the guard band. The upper panel shows the
background-only expected $\eps^2$ reach by mass and guard width. The lower panel makes
the mechanism explicit: wider guards remove nearby sideband information, increase
$\sigma_{A,\mathrm{ref}}$, and therefore soften the Gaussian-equivalent expected
reach.}
\label{fig:guard-band-reach-uncertainty}
\end{figure}

\begin{figure}[H]
\centering
\begin{subfigure}[t]{0.78\linewidth}
  \centering
  \graphicorplaceholder{\linewidth}{guard_band_figs/f_train_frac_vs_mass.png}
  \caption{Template fraction inside the training region.}
\end{subfigure}

\vspace{0.8em}

\begin{subfigure}[t]{0.78\linewidth}
  \centering
  \graphicorplaceholder{\linewidth}{guard_band_figs/reference_components_z0.png}
  \caption{Matched-reference and extraction-uncertainty components.}
\end{subfigure}
\caption{Leakage and reference-scale diagnostics used in the guard-band choice. The training-fraction panel shows the direct contamination pressure that disfavors too narrow a guard, while the reference-component panel checks that the matched background-only reference scale is consistent with the refit prescription used for the closure toys.}
\label{fig:guard-band-leakage-reference}
\end{figure}

\begin{figure}[H]
\centering
\begin{subfigure}[t]{0.49\linewidth}
  \centering
  \graphicorplaceholder{\linewidth}{guard_band_figs/eps2_recovery_vs_mass.png}
  \caption{Recovered over injected $\eps^2$.}
\end{subfigure}
\hfill
\begin{subfigure}[t]{0.49\linewidth}
  \centering
  \graphicorplaceholder{\linewidth}{guard_band_figs/delta_z_vs_mass.png}
  \caption{$Z$-calibration residuals.}
\end{subfigure}

\vspace{0.6em}

\begin{subfigure}[t]{0.49\linewidth}
  \centering
  \graphicorplaceholder{\linewidth}{guard_band_figs/pull_mean_vs_mass.png}
  \caption{Pull means.}
\end{subfigure}
\hfill
\begin{subfigure}[t]{0.49\linewidth}
  \centering
  \graphicorplaceholder{\linewidth}{guard_band_figs/pull_width_vs_mass.png}
  \caption{Pull widths.}
\end{subfigure}
\caption{Closure diagnostics entering the train-$2.25$ guard choice. Recovery and $Z$-calibration test the response to injected signals, while the pull panels diagnose amplitude bias and uncertainty calibration across the same mass and injected-strength grid. Curves show the toy-ensemble central values exported by the comparison suite; the uncertainty bands, where shown, are the toy-to-toy central spread for the same mass-strength slice rather than the uncertainty from one individual extraction. The train-$2.25$ configuration is well motivated because it suppresses training-region signal leakage while keeping the pull means near zero and the pull widths close to unity across the tested injections.}
\label{fig:guard-band-closure-diagnostics}
\end{figure}

\begin{figure}[H]
\centering
\begin{subfigure}[t]{0.98\linewidth}
  \centering
  \graphicorplaceholder{\linewidth}{guard_band_figs/eps2_recovery_box_z2.png}
  \caption{$Z_{\mathrm{inj}}=2$.}
\end{subfigure}

\vspace{0.8em}

\begin{subfigure}[t]{0.98\linewidth}
  \centering
  \graphicorplaceholder{\linewidth}{guard_band_figs/eps2_recovery_box_z5.png}
  \caption{$Z_{\mathrm{inj}}=5$.}
\end{subfigure}
\caption{Toy-level $\eps^2$ recovery distributions for all selected injected masses
(\SIlist{30;45;60;75;90;105;120}{MeV}). For each mass and training guard, the colored
box shows the central 50\% (interquartile) range of toy recoveries, the whisker bars
show the wider non-outlier toy range, and the black horizontal line inside the box
marks the toy median. The gray dashed horizontal reference lines mark 0.9, 1.0, and
1.1 recovery. The boxes show how the
narrow guard under-recovers injected signals, while train-$2.25$ restores the median
recovery much closer to unity before the wider-guard sensitivity cost becomes
dominant.}
\label{fig:guard-band-recovery-boxes}
\end{figure}

Once the training exclusion is fixed at $2.25\,\sigma_m$, there are two clean
blind-window geometries to compare. The older construction keeps the extraction likelihood at $1.64\,\sigma_m$ and treats the $1.64$--$2.25\,\sigma_m$ interval as a guard band. The newer construction uses a $2.25\,\sigma_m$ blind/extraction window, so the same bins are still excluded from training but are retained in the signal likelihood. The refmatched comparison holds the background-only reference convention fixed in both cases, so the remaining difference is the blind-window geometry itself. Figure~\ref{fig:refmatched-blind-geometry-reach} is the central reach comparison: it shows that the separate guard band is possible but unnecessary for the selected refmatched 2015 ensemble, because the widened blind window gives similar closure while recovering the signal tails that the guard-band construction discards. The older guard-band construction is kept as an appendix reference in Appendix~\ref{app:guard-band-construction}.

\begin{figure}[h!]
\centering
\graphicorplaceholder{0.82\linewidth}{guard_band_figs/refmatched_blind_geometries_eps2_95_reach.png}
\caption{Refmatched 2015 blind-geometry reach comparison. Both geometries exclude GP training data out to $2.25\,\sigma_m$ and define the injected strength with the matched background-only refit convention. The $2.25\,\sigma_m$ blind window has similar closure behavior but a lower expected $\eps^2_{95}$ proxy because the $1.64$--$2.25\,\sigma_m$ signal tails are included in the extraction likelihood instead of being discarded as a guard band.}
\label{fig:refmatched-blind-geometry-reach}
\end{figure}

\begin{figure}[h!]
\centering
\graphicorplaceholder{0.66\linewidth}{guard_band_figs/refmatched_blind_geometries_background_only_closure.png}
\caption{Background-only spurious-signal closure for the two refmatched blind
geometries. Both the guard-band construction and the $2.25\,\sigma_m$ blind/extraction window remain centered near zero pull in the tested 2015 functional-form ensemble.}
\label{fig:refmatched-blind-geometry-background}
\end{figure}

\begin{figure}[h!]
\centering
\graphicorplaceholder{0.82\linewidth}{guard_band_figs/refmatched_blind_geometry_pull_summary_z0135.png}
\caption{Compact pull-summary comparison for $Z_{\mathrm{inj}}=0,1,3,5$. The
$2.25\,\sigma_m$ blind window preserves the same signal-recovery behavior as the guard-band construction within the tested 2015 functional-form ensemble.}
\label{fig:refmatched-blind-geometry-pulls}
\end{figure}

\begin{figure}[h!]
\centering
\graphicorplaceholder{0.66\linewidth}{guard_band_figs/refmatched_blind_geometries_injected_amplitude_recovery.png}
\caption{Injected-amplitude recovery for the two refmatched blind geometries. The recovered amplitudes follow the injected reference scale similarly for the guard-band and widened-blind constructions.}
\label{fig:refmatched-blind-geometry-recovery}
\end{figure}

\begin{figure}[h!]
\centering
\graphicorplaceholder{0.66\linewidth}{guard_band_figs/refmatched_blind_geometries_z_calibration_residual.png}
\caption{$Z$-calibration residual comparison for the two refmatched blind geometries. In this 2015 ensemble the $2.25\,\sigma_m$ blind window and the $1.64\,\sigma_m$ blind plus $2.25\,\sigma_m$ guard construction have very similar closure behavior. The widened blind window is therefore selected because it keeps the same training protection while recovering the signal tails in the extraction likelihood and improving the expected reach.}
\label{fig:refmatched-blind-geometry-deltaz}
\end{figure}

\FloatBarrier

\subsection{Three toy-generation modes}

The validation program uses three complementary pseudoexperiment constructions.
Analytic functional-form pseudoexperiments provide the required smooth-background closure test because the true toy spectrum is independent of the GP model later used in the extraction.
Conditional blind-window GP toys are not used as a standalone closure requirement; they are a fast diagnostic of the local extraction likelihood when the sideband-trained GP prediction is held fixed.
Full procedural GP-mean/refit pseudoexperiments are the required closure stress test for signal absorption, because each toy is generated over the full mass range and then passed back through the same sideband-refit and extraction chain as the data.

\subsubsection{Analytic functional-form pseudoexperiments}

The first class of toys starts from an analytic fit to the smooth $e^+e^-$
invariant-mass spectrum and fluctuates pseudo-data around that closed-form background model. These studies isolate the question of gross spectral smoothness from the later question of GP conditioning. They are therefore useful as the first closure test before the full inference chain is exercised.

The primary analytic seeds are written explicitly here because their role is to define smooth positive toy means, not to claim a unique physical parameterization of the trident continuum. With $S(m;m_t,w)=[1+\exp\{-(m-m_t)/w\}]^{-1}$ and $u_d=(m-m_{\mathrm{mid},d})/m_{\mathrm{scale},d}$, the three selected dataset seeds are
\begin{align}
B_{2015}(m) &=
N_{15}\,S(m;m_{t,15},w_{15})\,(m-m_{0,15})^{a_{15}}
\nonumber\\[-0.3ex]
&\quad{}\times
\exp\!\left[-\frac{m-m_{0,15}}{\theta_{15}}\right]
\exp\!\left(c_{1,15}u_{15}+c_{2,15}u_{15}^2\right),
\label{eq:funcform-seed-2015}\\
B_{2016}(m) &=
N_{16}\,S(m;m_{t,16},w_{16})\,(m-m_{0,16})^{a_{16}}
\nonumber\\[-0.3ex]
&\quad{}\times
\exp\!\left[-\frac{m-m_{0,16}}{\theta_{16}}\right]
\exp\!\left(c_{1,16}u_{16}+c_{2,16}u_{16}^2\right),
\label{eq:funcform-seed-2016}\\
B_{2021}(m) &=
N_{21}\,S(m;m_{t,21},w_{21})\,m^{a_{21}}
\nonumber\\[-0.3ex]
&\quad{}\times
\exp\!\left[-\frac{m}{\theta_{21}}\right]
\exp\!\left(c_{1,21}m+c_{2,21}m^2\right).
\label{eq:funcform-seed-2021}
\end{align}
The 2015 and 2016 functions are independent fits of the same shifted sigmoid-power-exponential family with a normalized quadratic tail factor; they are set to zero below the fitted turn-on offset $m_{0,d}$ and are useful when the low-mass turn-on is better described by a shifted variable. The 2021 seed uses the unshifted sigmoid-power-exponential form with an exponential quadratic modifier, which was preferred for the 2021 1\% spectrum because the wider support and smoother low-mass turn-on did not require the shifted threshold parameter.

The fit selection algorithm is implemented in a deliberately conservative ROOT macro. For each dataset the macro builds a small list of positive candidate families, scans the configured lower edge of the fit interval, and for each candidate tries several seeded minimizations covering early turn-on, broad-tail, and late-turn-on shapes. Each trial is fit to the input histogram over the dataset-specific fit range with bin-integrated expectations, then evaluated with Pearson and Neyman goodness-of-fit on the selected support and with normalization checks on the full support range, scan range, and low/high sideband fractions. The primary seed is the first preferred candidate that has finite parameters, passes the support and sideband-fraction checks, and satisfies the configured Pearson target; if no preferred candidate passes, the best validated candidate is used. Here, support check means enough sideband bins after mass-dependent exclusion; sideband-fraction check means the retained sideband fraction; and Pearson target means the toy-validation $\chi^2$-like residual threshold. The selected function is then frozen, normalized to the observed number of events on the full toy-support range, and used to generate Poisson pseudo-data bin by bin from the normalized bin integrals.

The candidate families were chosen to span the few smooth structures that can mimic the observed spectra without introducing narrow features: a sigmoid threshold factor for turn-on behavior, a power-law rise, an exponential falloff for the high-mass tail, optional threshold shifts for datasets whose support begins near a sharp turn-on, and low-order exponential curvature for residual broad-shape freedom. The generalized-gamma and simpler sigmoid-power variants in the candidate overlay provide cross-checks that the chosen seed is not an artifact of one algebraic form. They are not used to inject additional model uncertainty in these closure plots; once the primary seed is selected, the closure question is whether the GP extraction recovers known signals from that fixed smooth truth.

Figure~\ref{fig:functional-form-seed-overview} shows the outcome of this construction
rather than an internal parameter dump. In each row, the left panel overlays the input
histogram with all enabled candidate families, while the right panel compares the
selected seed with the pseudoexperiment-ensemble mean. The 2015 shifted-tail family
follows both the turn-on and falling spectrum without introducing a narrow feature. The
2016 row shows the same family on the 10\% sample. The 2021 row selects the unshifted
\texttt{fSigPowExpQ} family, whose broad curvature follows the 1\% spectrum without
requiring a shifted threshold. Most importantly, each right-hand panel demonstrates
that the Poisson toy mean reproduces the selected analytic expectation and that the toy
normalization is set on the stated support rather than fixed inside the search window.

\begin{figure}[p]
\centering
\graphicorplaceholder{0.99\linewidth}{toy_generation_figs/funcform_publication_overview.png}
\caption{Functional-form seeds used to construct the analytic-background
pseudoexperiments. Each dataset row pairs the candidate-family comparison (left) with
the selected seed and the mean of the generated toy ensemble (right). The titles give
the family, the scan interval used for that generation campaign, and the full-support
normalization target. The 2015 row uses the shifted
sigmoid--power--exponential form with an additional smooth tail; the 2016 10\% row uses
an independently fitted member of the same family; and the 2021 1\% row uses the
unshifted \texttt{fSigPowExpQ} form. Shaded regions contribute to the normalization but
lie outside the displayed campaign's search interval. The 2016 42--210~MeV label is
retained to document the seed-generation study; the v3 result itself is restricted to
39--180~MeV. The corrected 2021 lower-bound study likewise uses 30--250~MeV model
support and searches 50--250~MeV, as documented below.}
\label{fig:functional-form-seed-overview}
\end{figure}

\paragraph{Injected-signal lower-bound scans.}
The lower-bound choice was made with signal-injection studies, not from the observed
limit curve. For every mass and $k_{\min}$ setting, the analytic background was
Poisson fluctuated, a full-range Gaussian signal was injected, the GP was refit to the
toy sidebands, and the blinded-window extraction was repeated. The 2015 and 2016 grids
used $Z_{\mathrm{inj}}=0,1,2,3,5$ at masses
\SIlist{45;60;75;90;105}{MeV} and \SIlist{60;90;105;120;150}{MeV}, respectively.
The corrected 2021 grid used $Z_{\mathrm{inj}}=0,1,3,5$ at
\SIlist{50;60;90;120;180;210}{MeV}. There are 100 toys per 2015 mass--strength cell
and 25 toys per 2016 and 2021 cell. The zero-injection cells test for a fitted false
signal; the nonzero cells test recovery, fitted uncertainty, and significance
calibration across the injected-strength range.

Table~\ref{tab:lslb-injection-pull-summary} makes the mean-pull information explicit.
The signed $Z_{\mathrm{inj}}=0$ entry is the equal-weight mean of the background-only
cell means over the tested masses. The signed $Z_{\mathrm{inj}}>0$ entry is the
equal-weight mean over all nonzero mass--strength cells. The RMS column is the root
mean square of those nonzero cell means and is sensitive to biases that change sign
across the grid; the last column is the median nonzero-injection pull width. No
uncertainty is assigned to the aggregate rows because the plots retain the
cell-by-cell toy-derived uncertainties used to judge whether a pattern is coherent.

\begin{table}[H]
\centering
\scriptsize
\begin{tabular}{lrrrrrr}
\toprule
Dataset & $k_{\min}$ & \makecell{Toys\\per cell} &
\makecell{Signed mean pull\\$Z_{\mathrm{inj}}=0$} &
\makecell{Signed mean pull\\$Z_{\mathrm{inj}}>0$} &
\makecell{RMS cell-mean pull\\$Z_{\mathrm{inj}}>0$} &
\makecell{Median pull width\\$Z_{\mathrm{inj}}>0$} \\
\midrule
2015 & 0.50 & 100 & $+0.235$ & $-1.449$ & 1.897 & 1.110 \\
     & 0.75 & 100 & $+0.146$ & $-0.508$ & 0.732 & 0.895 \\
     & 0.90 & 100 & $+0.090$ & $-0.259$ & 0.363 & 0.884 \\
     & $\mathbf{1.00}^{\ast}$ & 100 & $\mathbf{+0.029}$ & $\mathbf{-0.223}$ & $\mathbf{0.285}$ & $\mathbf{0.887}$ \\
     & 1.10 & 100 & $+0.037$ & $-0.161$ & 0.217 & 0.890 \\
\addlinespace
2016 10\% & $\mathbf{0.90}^{\ast}$ & 25 & $\mathbf{-0.016}$ & $\mathbf{-0.262}$ & $\mathbf{0.321}$ & $\mathbf{0.948}$ \\
     & 1.00 & 25 & $+0.178$ & $+0.036$ & 0.377 & 0.924 \\
     & 1.10 & 25 & $+0.222$ & $+0.120$ & 0.403 & 0.889 \\
\addlinespace
2021 1\% & 0.90 & 25 & $+0.061$ & $+0.041$ & 0.213 & 0.838 \\
     & 1.00 & 25 & $+0.016$ & $+0.004$ & 0.230 & 0.860 \\
     & $\mathbf{1.10}^{\ast}$ & 25 & $\mathbf{-0.005}$ & $\mathbf{-0.034}$ & $\mathbf{0.276}$ & $\mathbf{0.908}$ \\
\bottomrule
\end{tabular}
\caption{Aggregate pull diagnostics from the functional-form signal-injection scans.
An asterisk marks the lower-bound factor frozen for v3. Signed means and RMS values
refer to cell-level mean pulls, while $w_{\mathrm{pull}}$ is the fitted width of the
toy pull distribution in a mass--strength cell. The table should be read together with
the mass-resolved panels in Figures~\ref{fig:v3-lslb-closure-2015}--%
\ref{fig:v3-lslb-closure-2021}; a small aggregate mean alone can hide compensating
mass-dependent biases.}
\label{tab:lslb-injection-pull-summary}
\end{table}

Figures~\ref{fig:v3-lslb-closure-2015}--\ref{fig:v3-lslb-closure-2021} are the
physics-facing closure summaries for v3. Each panel uses analytic functional-form toys,
per-toy GP refitting, profiled background covariance, and toy-derived error bars. The
mean pull tests amplitude bias, the pull width tests the fitted uncertainty, the
$\hat Z-Z_{\mathrm{inj}}$ panel tests significance calibration when the reference scale
is matched, and the amplitude-bias panel tests signal recovery directly. These are
extraction-closure tests. Their reach panels and rankings are diagnostic proxies rather
than direct \CLs{} coverage measurements.

For 2015, the 100-toy study strongly rejects the overly flexible
$k_{\min}=0.5$ row. The 1.0 and 1.1 rows are close in pull-width calibration; the frozen
1.0 row has a background-only mean pull of $+0.029$, a signed nonzero-injection mean of
$-0.223$, and an RMS cell-mean pull of 0.285. It is retained as a conservative
production compromise even though the aggregate score slightly favors 1.1. For 2016,
the 25-toy scan selects $k_{\min}=0.9$ among 0.9, 1.0, and 1.1. In particular, the
background-only mean pull changes from $+0.222$ at 1.1 to $-0.016$ at 0.9, removing
the positive false-signal offset that motivated the follow-up. Across the injected
cells, 0.9 also reduces the RMS cell-mean pull from 0.403 to 0.321 and moves the median
pull width from 0.889 to 0.948. Its signed nonzero-injection mean is $-0.262$, so the
selection should be understood as the best joint calibration of the scanned rows, not
as exact zero bias in every injected cell.

The 2021 figure is not the older full-support study. It uses the regenerated
\path{final_1pct_invM.root} sample, an \texttt{fSigPowExpQ} truth model with
\SIrange{30}{250}{MeV} toy/data support, the \SIrange{50}{250}{MeV} search range, a
$2.25\,\sigma_m$ exclusion, $k_{\max}=9$, and a matched background-only refit reference.
All 1,800 requested toy rows have finite key fit quantities and successful refit and
$q_\mu$ statuses. The aggregate scores are nearly degenerate; $k_{\min}=1.1$ is frozen
because it gives the pull-width median closest to unity and the strongest extraction
proxy. Its background-only and nonzero-injection signed mean pulls are $-0.005$ and
$-0.034$, respectively; the larger 0.276 RMS reflects mass-to-mass variation that is
visible in Figure~\ref{fig:v3-lslb-closure-2021}. Pull widths below unity imply
conservative fitted uncertainties only if the toy ensemble faithfully represents the
production procedure. Direct coverage with the frozen upper-limit configuration
remains the final calibration requirement.

\begin{figure}[p]
\centering
\graphicorplaceholder{0.98\linewidth}{toy_generation_figs/v3_20260709_lslb_closure/fig40_style_2015_lslb_closure_suite_100toy.png}
\caption{2015 injected-signal closure as the RBF lower bound is varied. At each mass,
the plotted central value is the median of the nonzero
$Z_{\mathrm{inj}}=1,2,3,5$ cell summaries; each cell contains 100 full-refit
functional-form toys, and the error bars are derived from those toy ensembles. The
$k_{\min}=0.5$ behavior is incompatible with reliable signal recovery. The 1.0 and 1.1
curves form the stable neighborhood, and 1.0 is retained as the conservative v3
choice rather than as the minimum of the aggregate ranking score.}
\label{fig:v3-lslb-closure-2015}
\end{figure}

\begin{figure}[p]
\centering
\graphicorplaceholder{0.98\linewidth}{toy_generation_figs/v3_20260709_lslb_closure/fig40_style_2016_lslb_closure_suite_25toy.png}
\caption{2016 10\% injected-signal closure for
$k_{\min}=0.9,1.0,1.1$. Each point summarizes the nonzero
$Z_{\mathrm{inj}}=1,2,3,5$ cells at that mass, with 25 full-refit functional-form toys
per cell. The companion background-only cells have mean pulls $-0.016$, $+0.178$, and
$+0.222$ for 0.9, 1.0, and 1.1, respectively. Together with the smaller RMS
nonzero-injection mean pull and the pull width closer to unity, removal of the positive
background-only offset is the reason $k_{\min}=0.9$ is used for 2016 in v3.}
\label{fig:v3-lslb-closure-2016}
\end{figure}

\begin{figure}[p]
\centering
\graphicorplaceholder{0.98\linewidth}{toy_generation_figs/v3_20260709_lslb_closure/fig40_style_2021_search50_project30_refmatched_lslb_closure_suite_25toy.png}
\caption{Corrected 2021 1\% injected-signal closure for the
\texttt{fSigPowExpQ} toy truth. Each mass point summarizes the nonzero
$Z_{\mathrm{inj}}=1,3,5$ cells, with 25 toys per cell, \SIrange{30}{250}{MeV} model
support, and a \SIrange{50}{250}{MeV} search range. Defining the injected amplitude
with the matched background-only refit uncertainty removes the mixed-denominator
$\Delta Z$ offset seen in the superseded study. The 1.1 choice has the pull-width
median closest to unity and the strongest extraction-reach proxy; its signed mean pull
over the nonzero cells is $-0.034$. It is a calibration--reach choice, not the winner of
every scalar diagnostic.}
\label{fig:v3-lslb-closure-2021}
\end{figure}

\FloatBarrier

\subsubsection{Conditional blind-window GP toys}

The fastest toy mode keeps the sideband-trained GP fixed at a given test mass. The GP
fit is performed once on the sidebands, the blind-window mean vector and covariance are
constructed from that fit, a background realization is drawn, signal is added, and the
result is Poisson-sampled. This construction is ideal for high-statistics expected-band
studies and for testing the response of the extraction fit under the posterior
background model preferred by the real data.

\subsubsection{Full procedural GP-mean/refit pseudoexperiments}

The most realistic toy mode generates a full pseudoexperiment from the GP-derived
background expectation, injects signal with the full Gaussian template, and then reruns
the GP sideband-to-blind prediction on the toy before signal extraction. Here ``GP
mean'' means the full-range propagated count-space mean obtained from the sideband-fit
GP model of the observed spectrum. That propagated mean is used as the toy-truth
background source over the full histogram support; the toy is then fluctuated, signal
is added, and the ordinary mass-scan workflow refits the GP on the toy sidebands at
each scanned mass. This is the appropriate mode for diagnosing signal absorption,
because it directly tests how much of the injected signal may be absorbed by the
refitted background model or by the profiled nuisance treatment. The phrase
``refit-on-toy'' is therefore literal: the search stage does not reuse the original
data-conditioned blind-window prediction.

\subsubsection{Fixed-total GP-propagated toy-scan study}

The GP-propagated fixed-total toy mode remains a useful internal cross-check of the
scan machinery, but it is not the main closure evidence in this note. The main text
therefore emphasizes the clean year-matched functional-form scan suites above and the
refmatched full-refit closure studies below.

\subsection{Asimov calibration and background-only pseudoexperiments}

When toy strengths are specified in units of $\sigma_A$, the package first constructs a
reference extraction uncertainty $\sigma_{A,\mathrm{ref}}(m)$ by fitting the local
signal-strength model to a background-only reference spectrum,
\begin{equation}
\left(\hat A_{\mathrm{ref}},\sigma_{A,\mathrm{ref}}\right)
=
\mathcal{F}\!\left[\bm{n}_{\mathrm{ref}};\,\bm{b}(m),\bm{C}(m),\bm{w}(m)\right],
\end{equation}
where $\bm{b}(m)$ and $\bm{C}(m)$ are the GP-predicted blind-window mean and covariance,
and $\bm{w}(m)$ is the full-range Gaussian signal template restricted to the blinded
bins. The injected strength axis is then defined by
\begin{equation}
A_{\mathrm{inj}}(m) = Z_{\mathrm{inj}}\,\sigma_{A,\mathrm{ref}}(m).
\end{equation}

It is useful to state explicitly what $\sigma_{A,\mathrm{ref}}$ is \emph{not}. It is
not the uncertainty extracted from an individual toy after signal injection. Rather, it
is a calibration quantity computed once from a background-only reference problem at each
mass hypothesis. In the code this is exactly the object returned by evaluating the local
fit on the reference spectrum and reading off its fitted $\sigma_A$
\cite{Cowan1998,Cowan2011}. Its role is to define what is meant by an injected
``$1\sigma$,'' ``$2\sigma$,'' or ``$5\sigma$'' signal through
$A_{\mathrm{inj}}=Z_{\mathrm{inj}}\sigma_{A,\mathrm{ref}}$ before any toy-by-toy
fluctuations are generated.

In \textbf{Asimov mode}, the reference spectrum is deterministic,
\begin{equation}
n_{\mathrm{ref},i} = b_i,
\end{equation}
so $\sigma_{A,\mathrm{ref}}$ measures the nominal local sensitivity implied by the
fitted background model. In \textbf{Poisson mode}, the reference spectrum is first
fluctuated,
\begin{equation}
n_{\mathrm{ref},i} \sim \mathrm{Pois}\!\left(b_i\right),
\end{equation}
and the extracted $\sigma_{A,\mathrm{ref}}$ then inherits an extra layer of
background-only sampling noise. In both cases, the v3 headline closure studies use the
reference spectrum only to seed a background-only \emph{refit}. The reference
uncertainty is evaluated after imposing the same mass support, exclusion mask, kernel
bounds, and refit prescription used for the corresponding toy ensemble. This matched
background-only refit is the calibration reference used in the updated 2021 study and
in the year-specific lower-bound comparisons. Earlier studies that evaluated
$\sigma_{A,\mathrm{ref}}$ directly from a prefit Asimov spectrum are retained in
Appendix~\ref{app:prior-validation-results}; they do not define the v3 calibration.

This distinction matters for the later profile-likelihood interpretation. The Asimov
construction defines the expected sensitivity scale, while the background-only
pseudoexperiments are what test coverage, bias, and finite-sample departures from the
asymptotic approximation. Put differently: the Asimov dataset calibrates the injection
axis, but the pseudoexperiments remain the mechanism that profiles the background model
and measures how realistic background-only fluctuations propagate through the
extraction chain.

The later toy summaries therefore involve two related uncertainty scales. The reference
quantity $\sigma_{A,\mathrm{ref}}$ fixes the horizontal injection axis and the nominal
target significance $Z_{\mathrm{inj}}$. The toy-level $\sigma_A$ is returned by the fit
\emph{for that specific toy} after the background and signal fluctuations have been
realized. Pulls use the toy-level $\sigma_A$, whereas $Z_{\mathrm{inj}}$ is defined from
$\sigma_{A,\mathrm{ref}}$. Once the reference and toy fits are matched, the ensemble
pull and residual-significance calibrations should consequently be mutually
consistent, up to finite-sample and fit fluctuations. A large coherent separation is
a diagnostic of a reference-denominator or refit mismatch; it is not accepted as
normal behavior of an otherwise calibrated study.

\subsection{Poisson versus multinomial signal injection}

Signal injection is distinct from the choice of $\sigma_A$ reference. In
\textbf{Poisson-total injection}, each bin receives a Poisson draw with mean
$A_{\mathrm{inj}}w_i$, so the realized total signal fluctuates from toy to toy. In
\textbf{multinomial fixed-total injection}, the total number of injected events is
fixed and only the bin-to-bin allocation fluctuates. Since the present signal template
is normalized on the full histogram range, $A_{\mathrm{inj}}$ denotes the total local
signal yield and only the fraction $f_{\mathrm{blind}}A_{\mathrm{inj}}$ is expected to
appear inside the blinded bins for the selected $2.25\,\sigma_m$ window.

This distinction matters when interpreting uncertainty bars on validation plots. Only
Poisson-total injection produces a genuine horizontal spread in the realized injected
strength. Under fixed-total multinomial injection the horizontal coordinate is exact.

\subsection{Refit versus no-refit extraction displays}

The note now distinguishes two pseudoexperiment tests because they are answering
different questions. In the \textbf{no-refit} display, the GP sideband model
is fitted once on the real data and then held fixed while a background realization is
drawn in the blind window. This is the cleanest closure test of whether the extraction
fit can recover an injected signal when the background prediction is treated as the
correct local model. In the \textbf{full-refit} display, the toy is generated over the
full mass range and the GP is refit on the toy sidebands before the blind-window
extraction is repeated. That second test is slower and more aggressive, but it is also
the direct diagnostic for signal absorption by the refitted background model.

These tests should therefore not be judged by the same visual criterion. Agreement
between the no-refit pseudoexperiment display and the observed-data display is expected
and desirable because both are using the same sideband-conditioned background
description. The full-refit display is instead meant to stress the nuisance-model
freedom. When it visibly degrades the extracted signal, the correct interpretation is
not that the no-refit test was ``too easy,'' but that the full procedural toy is
probing how much the sideband refit can absorb signal-like structure.

\subsection{Representative single-mass extraction displays}

Representative extraction displays make the statistical procedure concrete at the level
of an individual mass hypothesis. The 2015 \SI{30}{MeV} display and the 2016 10\%
\SI{58}{MeV} display are retained because they show, in real data, how the profiled
background deformation competes with the narrow signal template. The 2016 example is
particularly useful because the best-fit amplitude is negative.
Figure~\ref{fig:data-fit-examples} is included to make those competing fit components
visible before moving to ensemble summaries.

\begin{figure}[h!]
\centering
\begin{subfigure}[t]{0.88\linewidth}
  \centering
  \graphicorplaceholder{\linewidth}{fit_example_figs/2015/m030MeV_blind_fit.png}
  \caption{2015 data at \SI{30}{MeV}.}
\end{subfigure}

\vspace{1.0em}

\begin{subfigure}[t]{0.88\linewidth}
  \centering
  \graphicorplaceholder{\linewidth}{fit_example_figs/2016/m058MeV_blind_fit.png}
  \caption{2016 10\% data at \SI{58}{MeV}, including a negative $\hat A$.}
\end{subfigure}
\caption{Representative blind-window fits using data. These displays show the fit
components directly and serve as concrete examples of the local extraction machinery
before one moves to ensemble diagnostics.}
\label{fig:data-fit-examples}
\end{figure}

\subsection{Pull, coverage, and significance-calibration diagnostics}

The central ensemble-level validation quantities are the pull,
\begin{equation}
\mathrm{pull} = \frac{\hat A - A_{\mathrm{inj}}}{\sigma_A},
\end{equation}
its mean and width as a function of mass and injected strength, and the empirical
coverage of the nominal one- and two-standard-deviation intervals. These are standard
closure diagnostics in HEP statistical analyses because they test, respectively,
whether the estimator is biased, whether the quoted uncertainty is correctly calibrated,
and whether the reported intervals achieve the claimed frequentist coverage
\cite{Cowan1998,Cowan2011}.

The main closure evidence in this section comes from the functional-form
pseudoexperiments, where
the smooth-spectrum truth is an analytic toy seed rather than a GP-generated background
ensemble. Earlier GP-generated pull, coverage, and $\Delta Z$ displays remain useful
internal checks of the extraction layer, but they are not shown as headline validation
figures because the refmatched functional-form pull and calibration plots are the more
direct closure tests.

\subsubsection{Representative pull distributions}

The year-specific v3 figures above summarize closure across mass and injected strength.
Figure~\ref{fig:refmatched-train225-pulls-2015} complements those summaries by showing
the toy-level pull distributions at one representative 2015 mass. It uses the
reference-matched definition given above: the zero-injection panel tests for a spurious
signal, while the nonzero panels test recovery after the same sideband-refit procedure
has been applied to the toys.

\begin{figure}[p]
\centering
\begin{subfigure}[t]{0.60\linewidth}
  \centering
  \graphicorplaceholder{\linewidth}{injection_extraction_figs/refmatched_train2p25_pull_hists_2015/pull_hist_2015_m075MeV_Z0.png}
  \caption{$Z_{\mathrm{inj}}=0$.}
\end{subfigure}
\par\vspace{0.25em}
\begin{subfigure}[t]{0.60\linewidth}
  \centering
  \graphicorplaceholder{\linewidth}{injection_extraction_figs/refmatched_train2p25_pull_hists_2015/pull_hist_2015_m075MeV_Z2.png}
  \caption{$Z_{\mathrm{inj}}=2$.}
\end{subfigure}
\par\vspace{0.25em}
\begin{subfigure}[t]{0.60\linewidth}
  \centering
  \graphicorplaceholder{\linewidth}{injection_extraction_figs/refmatched_train2p25_pull_hists_2015/pull_hist_2015_m075MeV_Z5.png}
  \caption{$Z_{\mathrm{inj}}=5$.}
\end{subfigure}
\caption{Representative 2015 train-$2.25$ refmatched pull distributions at
\SI{75}{MeV}. The zero-injection panel is the direct spurious-signal test, while the
$Z_{\mathrm{inj}}=2$ and $Z_{\mathrm{inj}}=5$ panels show the injected-signal pull
response under the same full-refit widened-blind prescription.}
\label{fig:refmatched-train225-pulls-2015}
\end{figure}

The earlier two-panel dataset summaries are superseded by
Figures~\ref{fig:v3-lslb-closure-2015}--\ref{fig:v3-lslb-closure-2021}. They are retained
in Appendix~\ref{app:prior-validation-results} because they document the development of
the matched-reference procedure, but they do not use the selected three-dataset lower-bound
choices and should not be read as the v3 headline closure evidence.

\FloatBarrier

The earlier fixed-background/no-refit functional-form plots remain useful as internal
diagnostics because they isolate the local likelihood once the sideband-conditioned GP
prediction is held fixed. They are not the main closure evidence in this note and are
therefore not shown as headline figures here. The full-refit panels above are the
appropriate studies for assessing signal absorption by the sideband refit.

\subsection{Auxiliary combined-study displays}

The earlier combined-search-power scans and illustrative no-refit/full-refit extraction
displays have been moved to Appendix~\ref{app:prior-validation-results}. They remain
useful demonstrations of how a common $\eps^2$ maps into dataset-dependent signal
yields, but their provisional injections are not direct coverage or closure evidence
for the frozen v3 result. The physics-facing simultaneous limit and its matching
observed-equivalent full-sample projection are presented in
Section~\ref{sec:results-combined-151621-1pct}.
